\documentclass{article}\begin{document}For an indexed family of actual source and destination functors and component cleavages, choice builds a cleavage on the source leg of the single coproduct apex. TagCartesian proves universal factorization against every arrow in the apex: a morphism into a tagged target must have that tag, and the component factorization is unique. WeightedChoice then uses the actual destination leg and the contravariant presheaf pull, not transport across an equivalence. On global weighted strict fibres lift\_tag computes the entire chosen arrow, and map\_tag proves the image of the component's vertical reindexing map equals the single-span vertical reindexing map by cancelling against that chosen cartesian lift. The resulting strictSquare is an equality on all vertical arrows. The two strict inverse fibre presentations give presentedSquare. compareTag identifies these presented inclusions with restrictions of the previous actual compare functor via composition of endpoint-aligned total maps; forward assembles their base squares. For matching base, grade and payload universe levels, All.parts picks the identity-span cleavage at zero and the original positive-power cleavages otherwise. All.evaluationBase is therefore the actual isomorphism from first pulling in the weighted single span and then evaluating, to first evaluating then pulling in the tagged iterates, for each base arrow and fixed payload. It uses the chosen inverse equivalence's mate, not an identity replacement. This is an individual natural isomorphism of functors; base identity/composition diagrams for these new squares, their mixed compatibility with \#197's payload square, and bundled pseudonaturality are still required. The concrete all-length size restriction has not been relaxed. No flattening, monad or algebra recovery. Manuscript 458ff36a lines 929--953; reconstruct the comparison by reading the factorization in NativeCleavage.Reindex, the two stages of CoproductSpan.compare, then the use of SumChange.inverseSquare and taggedBase.\end{document}
